%% ma: L12-B14-0136
%% lop: 12
%% chuong: 5
%% bai: 14
%% dang: 12.14.22
%% loai: TN4
%% muc: VD
%% dap_an: "C"
%% nguon: "Ảnh giáo viên gửi 25/11/2026 (tài liệu ôn tập Phương trình mặt phẳng), Phần I Câu 43"
%% kien_thuc: "Bài 14 (phương trình mặt phẳng theo đoạn chắn, điều kiện hai mặt phẳng vuông góc)"
%% trang_thai: chinh-thuc
\begin{ex}
Trong không gian $Oxyz$, cho $3$ điểm $A(1;0;0)$, $B(0;b;0)$, $C(0;0;c)$ trong đó $bc\ne0$ và mặt phẳng $(P)\colon y-z+1=0$. Mối liên hệ giữa $b,c$ để mặt phẳng $(ABC)$ vuông góc với mặt phẳng $(P)$ là
\choice{$2b=c$}{$b=2c$}{\True $b=c$}{$b=3c$}
\loigiai{$(ABC)\colon x+\dfrac yb+\dfrac zc=1$ có vectơ pháp tuyến $\vec n_1=\left(1;\dfrac1b;\dfrac1c\right)$; $(P)$ có $\vec n_2=(0;1;-1)$. $(ABC)\perp(P)\Leftrightarrow\vec n_1\cdot\vec n_2=0\Leftrightarrow\dfrac1b-\dfrac1c=0\Leftrightarrow b=c$.}
\end{ex}
